\documentclass[10pt,a4paper]{amsart}
\usepackage[margin=19mm]{geometry}
\usepackage{amsmath,amssymb,mathtools}
\usepackage[scr=rsfs,cal=euler]{mathalfa}
\usepackage{xfrac}
\usepackage[unicode,hidelinks,bookmarks=false]{hyperref}
\newcommand{\ie}{i.e.\ }
\newcommand{\cf}{cf.\ }
\newcommand{\resp}{respectively\ }
\newcommand{\Subsection}{\subsection}
\providecommand{\nobreakdash}{\nobreak\hbox{-}}
\setlength{\emergencystretch}{2em}
\multlinegap0pt
\usepackage{mathtools}
\usepackage{amssymb}
\usepackage{bm}
\usepackage[normalem]{ulem}
\usepackage[shortlabels]{enumitem}
\newtheorem{proposition}{Proposition}
\newtheorem{lemma}{Lemma}
\title{The Planar Coleman--Gurtin Model with Beltrami Conductivity}
\author{Francesco Di Plinio}
\date{Source: arXiv:2603.05983v1 (CC BY 4.0)}
\begin{document}
\maketitle

\noindent\emph{Source and licence.} The Planar Coleman--Gurtin Model with Beltrami Conductivity. Version arXiv:2603.05983v1 (CC BY 4.0), distributed by arXiv under the Creative Commons Attribution 4.0 International licence, http://creativecommons.org/licenses/by/4.0/, as recorded in the arXiv licence metadata for this exact version and shown on the abstract page https://arxiv.org/abs/2603.05983v1. Full licence notice: \texttt{licenses/arxiv-2603.05983-CC-BY-4.0.txt}.\par\smallskip

\noindent\textit{Editorial scope.} Counted unit: the article's lemma lem:phi-lip-H2 (source0.tex lines 1740--1745). The article's complete original proof of it is delivered with it (source0.tex lines 1747--1823, one \texttt{proof} environment of 77 lines). No mathematics is written, repaired or re-derived: the statement and the proof are the article's own lines, in the article's own notation, and no retained line is substituted or re-flowed.  Retained uncounted as the source-local context the proof needs: lines 224--227; lines 262--264; lines 265--267; lines 714--726.  Nothing was omitted from any retained span, so the package carries no omission record.  No retained line contains a citation, so the wrapper carries no bibliography and no bibliography entry.  No retained line contains a figure environment, an image-inclusion command or drawing code, so no image is delivered and the package declares no asset dependency.  Editorial formatting only, in addition: an AMS \texttt{amsart} preamble that restates the article's own declarations and macros used by the retained lines, namely the environments \texttt{proposition}, \texttt{lemma} on separate counters, together with the standard packages those lines need.  The retained environments print in this document as Proposition 1, Lemma 1 in order of appearance, whereas the article prints the same items with the numbering of its own full document.  The complete native source of the article as distributed in the arXiv e-print for this exact version is bundled unchanged as \texttt{sources/195830/source0.tex}.
\begin{equation}\label{eq:ellipticity}
m(k)|\xi|^{2}\le \sigma_\mu(x)\xi\cdot\xi\le M(k)|\xi|^{2}
\qquad\text{ a.e. }x\in\Omega,\, \forall \xi\in\mathbb R^{2}.
\end{equation}

The dilatation $\mu:\Omega \to \mathbb C$ satisfies the ellipticity condition \begin{equation} \label{e:ellipticitymu}
	\|\mu\|_{L^\infty(\Omega)} \le k < 1.
\end{equation}Assumption \eqref{e:ellipticitymu} is standing throughout the article. The regularity hypothesis

 \begin{equation} \label{e:regmu}
\mu \in W^{1,p_0}(\Omega), \qquad p_0\geq 2
\end{equation}

\begin{proposition} \label{prop:W2p-A}
Assume \eqref{e:ellipticitymu} holds. Also assume that \eqref{e:regmu} holds for some $p_0\geq 2.$ 
Let $p \in (1,p_0)$, $g\in L^{p}(\Omega)$ and let $u\in H^{1}_{0}(\Omega)$ be the (unique) weak solution of
\[
A_\mu u = g \qquad\text{in }H^{-1}(\Omega).
\]
Then $u\in W^{2,p}(\Omega)\cap H^{1}_{0}(\Omega)$ and
\begin{equation}\label{eq:W2p-est-A}
\|u\|_{W^{2,p}(\Omega)}
\le C\left(\|g\|_{L^{p}(\Omega)}+\|u\|_{L^{2}(\Omega)}\right).
\end{equation}
The constant $C$ depends on $p$, $\Omega$, $k$ and $\|\mu\|_{W^{1,p_0}(\Omega)}$. 
\end{proposition}

\begin{lemma}[Local Lipschitz of $\phi$ in ${V^2}$]\label{lem:phi-lip-H2}
Assume   \eqref{e:regmu} holds for some $p_0 \geq 2$. Then 
\begin{equation}\label{eq:phi-lip-H2}
\|\phi(u_2)-\phi(u_1)\|_{{V^2}}\le \mathcal I(\|u_1\|_{{V^2}}+ \|u_2\|_{{V^2}}) \|u_2-u_1\|_{{V^2}}.
\end{equation}
\end{lemma}

\begin{proof}
Set $w\coloneqq u_2-u_1$, and fix $p\in (\frac43,p_0)$. As before, Sobolev embeddings and  the estimates of Proposition~\ref{prop:W2p-A}   
yield the estimate 
\begin{equation}\label{eq:Linfty-L4-sec8}
\|u\|_{\infty}+\|\nabla u\|_{L^4(\Omega)}\le C\|u\|_{W^{2,p}}\le C\|u\|_{{V^2}}
\end{equation}
Using \textbf{(P2)}, and \eqref{eq:Linfty-L4-sec8} we then have 
\begin{equation}\label{eq:phi-derivs-bdd-sec8}
\|\phi^{(1+j)}(\theta)\|_{\infty}\le \mathcal I(\|u_1\|_{{V^2}}+ \|u_2\|_{{V^2}}), \qquad j=0,1,2\end{equation}
for any intermediate function $\theta$ between $u_1$ and $u_2$.
The chain rule reads
\begin{equation}\label{eq:chain-Amu-sec8}
A_\mu(\phi(u))=\phi'(u)A_\mu u-\phi''(u)(\sigma_\mu\nabla u)\cdot\nabla u
\qquad\text{in }L^2(\Omega).
\end{equation}
Applying \eqref{eq:chain-Amu-sec8} to $u_2$ and $u_1$ and subtracting gives
\[
A_\mu(\phi(u_2)-\phi(u_1))=\mathrm{I}+\mathrm{II}-\mathrm{III},
\]
where
\[
\begin{split}&
\mathrm{I}\coloneqq \phi'(u_2)A_\mu w,
\qquad
\mathrm{II}\coloneqq \left(\phi'(u_2)-\phi'(u_1)\right)A_\mu u_1,
\\
&
\mathrm{III}\coloneqq
\phi''(u_2)(\sigma_\mu\nabla u_2)\cdot\nabla u_2
-\phi''(u_1)(\sigma_\mu\nabla u_1)\cdot\nabla u_1.
\end{split}
\]
We estimate these terms in $L^2$.

\smallskip
\noindent\emph{Term $\mathrm{I}$.} Using \eqref{eq:phi-derivs-bdd-sec8},
\[
\|\mathrm{I}\|_{V^0}\le \|\phi'(u_2)\|_{\infty}\|A_\mu w\|_{V^0}
\le  \mathcal I(\|u_1\|_{{V^2}}+ \|u_2\|_{{V^2}})\|w\|_{{V^2}}.
\]
\smallskip
\noindent\emph{Term $\mathrm{II}$.} By the mean value theorem,
$\phi'(u_2)-\phi'(u_1)=\phi''(\theta)w$ for an intermediate function $\theta$, hence a combination of \eqref{eq:Linfty-L4-sec8} and \eqref{eq:phi-derivs-bdd-sec8} entails
\[
\|\mathrm{II}\|_{V^0}
\le \|\phi''(\theta)\|_{\infty}\|w\|_{\infty}\|A_\mu u_1\|_{V^0}
\le \mathcal I (\|u_1\|_{{V^2}} + \|u_2\|_{{V^2}})\|w\|_{{V^2}}.
\]
\smallskip
\noindent\emph{Term $\mathrm{III}$.} Rewrite
\[
\mathrm{III}=
\phi''(u_2)\left(Q_2-Q_1\right)+\left(\phi''(u_2)-\phi''(u_1)\right)Q_1,
\qquad
Q_j\coloneqq (\sigma_\mu\nabla u_j)\cdot\nabla u_j, \quad j=1,2.
\]
For the first part, note that
\[
Q_2-Q_1=(\sigma_\mu\nabla w)\cdot\nabla u_2+(\sigma_\mu\nabla u_1)\cdot\nabla w,
\]
so, by ellipticity \eqref{eq:ellipticity} and H\"older's inequality
\[
\|\phi''(u_2)(Q_2-Q_1)\|_{V^0}
\le \|\phi''(u_2)\|_\infty\|\nabla w\|_{L^4(\Omega)}\left(\|\nabla u_2\|_{L^4(\Omega)}+\|\nabla u_1\|_{L^4(\Omega)}\right)
\le \mathcal I(\|u_1\|_{{V^2}}+ \|u_2\|_{{V^2}})\|w\|_{{V^2}},
\]
having relied upon \eqref{eq:Linfty-L4-sec8} and \eqref{eq:phi-derivs-bdd-sec8}.
For the second part, again by the mean value theorem,
$\phi''(u_2)-\phi''(u_1)=\phi'''(\zeta)w$ for an intermediate $\zeta$, hence
\[
\|(\phi''(u_2)-\phi''(u_1))Q_1\|_{V^0}
\le \|\phi'''(\theta)\|_\infty\|w\|_{\infty}\|Q_1\|_{V^0}.
\]
Finally, $\|Q_1\|_{V^0}\leq C \|\nabla u_1\|_{L^4(\Omega)}^2\le  \mathcal I(\|u_1\|_{{V^2}} )$ by \eqref{eq:Linfty-L4-sec8},
and it follows that the last display is $ \leq \mathcal I(\|u_1\|_{{V^2}}+ \|u_2\|_{{V^2}})\|w\|_{{V^2}}$ as well.
Collecting the estimates for $\mathrm{I},\mathrm{II},\mathrm{III}$ finally yields \eqref{eq:phi-lip-H2}.
\end{proof}

\end{document}
